\chapter{Conclusions and Future Work}
\section{Conclusion}
In this thesis a modular ionospheric model was presented to provide coupling to a global magnetohydrodynamic model of the magnetosphere. Empirical models were implemented for the ionospheric conductivity with future opportunity to provide multiple new models. Additionally, a spherical discretization of the ionospheric surface equation was provided with a physically grounded boundary condition. The application shown provides an extensible framework for ionosphere simulation with the ability to provide output across different conductivity models, equations and boundary conditions. Through the solver and preconditioner parameter sweep, substantial performance improvements in the runtime of the application were obtained, reducing solve time from approximately 140 seconds to 10.

The qualitative analysis of the simulation output provided verification of the general structure of the produced potentials and the boundary condition behaviour. While there is still work remaining on certain portions of the application such as the magnetosphere-ionosphere coupling process, the model acts as a proof of concept for an extensible ionosphere simulation. Through additional verification and optimization as detailed in the next section, this model will provide the dissipative boundary condition that will allow for accurate space weather simulation.
\section{Future Work}
\subsection{Magnetosphere-Ionosphere Coupling}
Magnetosphere-ionosphere coupling is the immediate priority for future work, enabling this model to be used for practical space-weather prediction purposes. The incompleteness of the magnetosphere-ionosphere coupling has limited the analysis that is possible to perform on the output of the model. Once the solver loop described in Section~\ref{sec:application_structure} can be closed by providing time-stepped radial current data from the simulation, new strategies can be investigated to provide potentially increased solver performance. Additionally, this will allow for validation of the model's correctness over a range of IMF source terms.
\subsection{Verification}
A comprehensive verification of the predictions of the model is planned once the magnetosphere-ionosphere coupling is complete. This will involve testing the model over a wide range of IMF source terms and validating its output against reference potentials. Possible targets for validation include techniques such as the AMIE procedure~\citep{richmond1988amie} and LOMPE~\citep{local_laundal_2022}, which provide procedures to map high-latitude electric fields from sets of localized observational data, and electrostatic potential distributions produced by the SuperDARN network through the approach described in~\citet{ruohoniemi1998superdarn}.

A more in-depth analysis of the correctness of the results produced by this model's boundary condition is another priority. The polar boundary condition provided qualitatively good results. However, additional validation of the potentials produced by this boundary condition is still planned. This will likely involve testing the boundary condition against a series of manufactured solutions to verify its correctness.

\subsection{Solver and Performance Extensions}
The analysis in Section~\ref{sec:convergence_study} of this report provided substantial gains in performance. Iteration count was reduced from approximately 1500 to 100, with an improvement in solve time from approximately 140 to 10 seconds. However, the goal for the final model is to achieve sub-second solve times under most conditions. To achieve this, a more in-depth study of the convergence behaviour of the simulation is required. This would include finer parameter tuning, testing on a more realistic system with dedicated cores or a networked MPI cluster and a broader study of different solvers and preconditioning methods.

Avenues to be explored are optimized MPI rank data distribution including testing interleaved row storage methods and using the Trilinos package Zoltan to automatically partition data. Additionally, experimenting with the \texttt{GCRO-DR} GMRES solver that recycles Krylov subspace information across solves is planned. The last optimization is particularly promising as it is designed for systems where the matrix structure is constant with sequentially changing source vectors. Depending on the time scale required by the implementation, it could be used to recycle data between solves potentially further improving performance~\citep{parks2006}.
\subsection{Model Extensions}
One of the primary design choices in this application was the extensibility provided by the application structure. In order to leverage this, future plans include additional empirical models for the conductivity, implementing a new conductivity tensor and a new boundary condition. Alternate empirical models for the Hall and Pedersen conductivities are available as shown in~\citet{cousins2015} and~\citet{local_laundal_2022}. Additionally, the implementation of a plasma-parameter approach utilizing the NRLMSIS~\citep{nrlmsis_emmert_2020} and IRI~\citep{international_bilitza_2022} models in conjunction with the equations described in Equation~\ref{eq:physical_conduct} is planned. This would provide an alternative source for conductivity data while also providing values for the parallel conductivity. The alternative conductivity tensor proposed by~\citet{goodman1995} is also planned to be implemented as this would remove the requirement for the parallel conductivity entirely, and would allow the comparison of the two conductivity tensors. An alternative boundary condition, that is in the process of implementation, utilizing a five-point stencil mapped onto a ring of ghost points around the pole is also going to be added to the model. This is likely to result in better convergence behaviour than the divergence theorem approach, but does not provide continuity over the pole as a constraint. Together these planned extensions demonstrate how the modular structure of the application allows for exploration of multiple physical models without changes to the solver.

